\section{Future Directions}
\label{sec:future}

We identify five directions likely to change how spatial resolution enhancement is approached computationally. These directions are not speculative; they reflect trends already underway in the field that will mature over the next three to five years.

Foundation models for spatial omics represent the most transformative opportunity. Large-scale pre-trained models such as Nicheformer, Novae, OmiCLIP, and spatial extensions of scGPT \citep{cui2024scgpt} and CellPLM are learning transferable representations of cellular identity and spatial context from atlas-scale data. These models could reduce dependence on per-dataset references by providing universal cell-type embeddings that generalise across tissues and platforms. Whether representations learned primarily from dissociated single-cell data transfer to the spatial domain, where neighbourhood context carries biological meaning absent in suspension, remains an open question. Nicheformer explicitly encodes spatial neighbourhood composition during pre-training, making it the most directly relevant to deconvolution and reconstruction tasks. scGPT, by contrast, was designed for dissociated cells and requires fine-tuning on spatial data to capture tissue architecture.

The Visium HD era and sub-cellular sequencing will reframe rather than eliminate the resolution problem. Visium HD's 2~\textmu m bins approach single-cell resolution for many tissue types, potentially rendering traditional spot-level deconvolution unnecessary. Yet computational challenges persist: at 2~\textmu m, individual bins capture very few UMIs, requiring aggregation or imputation to achieve reliable expression estimates. The deconvolution problem shifts from ``which cell types are in this spot?'' to ``how should we aggregate sparse bins into biologically meaningful units?'' Furthermore, the vast installed base of standard Visium data ensures that computational resolution enhancement will remain relevant for re-analysis of existing datasets for years to come.

Reference-free and self-supervised methods offer a path toward resolution enhancement without external references. Current reference-free approaches (STdeconvolve, BayesTME) are limited to Stage~I and produce coarse outputs. Self-supervised learning, training on the spatial data itself using spatial consistency, image features, or multi-scale structure as supervision signals, could enable reference-free methods that achieve Stage~II or even Stage~III resolution. Contrastive learning frameworks that enforce consistency between a spot's expression and its histological appearance are a promising direction.

Standardised benchmarks and community evaluation are urgently needed. The field needs shared benchmark resources analogous to the Open Problems in Single-Cell Analysis initiative. Key requirements include multi-platform datasets with matched sequencing-based and imaging-based ST on the same tissue, agreed-upon metrics that capture both accuracy and calibration, standardised preprocessing to eliminate confounding differences between evaluation studies, and living leaderboards that track method performance as new tools are published.

Multi-omics spatial integration will expand the scope of resolution enhancement. Emerging technologies measure not only transcriptomes but also proteins (CODEX, IMC), chromatin accessibility (spatial ATAC-seq), and metabolites in spatial context. Integrating these modalities computationally, for example using spatial proteomics to anchor transcriptomic deconvolution or using chromatin state to inform cell-type identity, will require new multi-modal frameworks that go beyond the expression-centric methods reviewed here. The challenge is not merely technical but also conceptual: defining what ``single-cell resolution'' means when different modalities have different native resolutions and measurement characteristics.
